\documentclass[12pt]{ctexart}

% ============ 页面尺寸 ============
\usepackage[paperwidth=210mm, paperheight=297mm,
top=2cm, bottom=1.5cm,
left=1.8cm, right=1.8cm,
headheight=15pt]{geometry}

% ============ 必要宏包 ============
\usepackage{xcolor}
\usepackage{tcolorbox}
\usepackage{tikz}
\usetikzlibrary{arrows.meta}
\usepackage{fancyhdr}
\usepackage{calc}
\usepackage{graphicx}

\tcbuselibrary{breakable}

% ============ 页眉页脚 ============
\pagestyle{fancy}
\fancyhf{}
\fancyhead[L]{\sffamily 错题收录}
\fancyfoot[C]{\sffamily\color{black}\small 第\,\thepage\,页}
\renewcommand{\headrulewidth}{1pt}

% ============ 间距控制 ============
\newlength{\problemTopSkip}
\newlength{\problemBottomSkip}
\setlength{\problemTopSkip}{0.02\paperheight}  % 页首基础空白
\setlength{\problemBottomSkip}{6mm}            % 只留基础间距，大留白交给 \vfill

% ============ 计数器 + 公共换页钩子 ============
\newcounter{problemcount}

% 每题结束：本页题数 +1；用 \vfill 吃掉剩余高度（自动均分）；
% 满 3 题强制换页并把计数器清零。
\newcommand{\problempagebreak}{%
	\stepcounter{problemcount}%
	\par\vfill
	\ifnum\value{problemcount}=3\relax
	\setcounter{problemcount}{0}%
	\clearpage
	\fi
}

% ============ 环境 1：LaTeX 画图版（题干在上，图在下方居中） ============
\newtcolorbox{problem}[2]{%
	frame empty,
	colback=white,
	boxrule = 0pt,
	left    = 0pt,
	right   = 0pt,
	top     = 0pt,
	bottom  = 0pt,
	breakable = true,
	before skip = \problemTopSkip,
	after skip  = \problemBottomSkip,
	before upper = {\noindent\textbf{#1\quad #2}\par\smallskip},
	after = {\problempagebreak},
}

% ============ 环境 2：外部图片右对齐版 ============
% 用法：\begin{problemfig}{题号}{题干}{图片路径}{右侧图片宽度比例}
	%   #1 = 题号    #2 = 题干
	%   #3 = 图片文件名（相对/绝对路径均可）
	%   #4 = 右侧 minipage 宽度占 \linewidth 的比例，建议 0.30 ~ 0.40
	%
	% 左侧题干固定占 0.55\linewidth，右侧图片占 #4\linewidth，
	% 中间 \hfill 弹性填充；图片宽高双限制（不超过 0.25\textheight），
	% keepaspectratio 保证不变形。
	\newtcolorbox{problemfig}[4]{%
		frame empty,
		colback=white,
		boxrule = 0pt,
		left    = 0pt,
		right   = 0pt,
		top     = 0pt,
		bottom  = 0pt,
		breakable = true,
		before skip = \problemTopSkip,
		after skip  = \problemBottomSkip,
		% ---- 左侧：题干 ----
		before upper = {%
			\begin{minipage}[t]{0.55\linewidth}%
				\vspace{0pt}%           % 修正 [t] 顶部对齐基线
				\noindent\textbf{#1\quad #2}\par\smallskip
			},
			% ---- 右侧：图片 ----
			after upper = {%
			\end{minipage}%
			\hfill
			\begin{minipage}[t]{#4\linewidth}%
				\vspace{0pt}%
				\centering
				\includegraphics[width=\linewidth,
				height=0.25\textheight,
				keepaspectratio]{#3}%
			\end{minipage}%
		},
		after = {\problempagebreak},
	}
	
	\begin{document}
		\setlength{\parindent}{0pt}
		
		% ============================================================
		% 示例 1：LaTeX 画图版
		% ============================================================
		\begin{problem}{01}{台阶周长}
			如图所示，是一个向右下方延伸的3级台阶平面图。
			已知台阶的总高度为 10\,cm，总长度为 15\,cm。
			这个图形的周长是 \underline{\hspace{2cm}}\,cm.
			
			\vspace{1ex}
			\begin{center}
				\resizebox{0.75\linewidth}{!}{%
					\begin{tikzpicture}[>=Stealth]
						\draw[thick] (0,0) -- (5,0) -- (5,3.33) -- (10,3.33) -- (10,6.66) -- (15,6.66) -- (15,10) -- (0,10) -- cycle;
						\draw[thick, dashed, red] (0,0) rectangle (15,10);
						\draw[<->] (0,-0.7) -- (15,-0.7) node[midway, below, font=\LARGE\bfseries] {15cm};
						\draw[<->] (-0.7,0) -- (-0.7,10) node[midway, left, font=\LARGE\bfseries] {10cm};
					\end{tikzpicture}%
				}
			\end{center}
		\end{problem}
		
		% ============================================================
		% 示例 2：外部图片右对齐版
		%   请把 figures/step.png 换成你自己的图片文件
		% ============================================================
		\begin{problemfig}{02}{台阶周长（外部图片）}{Picture/三年级暑假数学第九题.png}{0.38}
			如图所示，是一个向右下方延伸的3级台阶平面图。
			已知台阶的总高度为 10\,cm，总长度为 15\,cm。
			这个图形的周长是 \underline{\hspace{2cm}}\,cm.
		\end{problemfig}
		
		% ============================================================
		% 示例 3：混用，验证每页 3 题 + 弹性留白
		% ============================================================
		\begin{problemfig}{03}{三角形面积}{Picture/三年级暑假数学第九题.png}{0.35}
			下图中三角形的底为 8\,cm，高为 5\,cm，求其面积。
		\end{problemfig}
		
		\begin{problem}{04}{下一个图形题}
			第 4 题，应出现在第 2 页页首。
		\end{problem}
		
		\begin{problem}{05}{再下一个图形题}
			第 5 题。
		\end{problem}
		
	\end{document}